\SourceLeaf{190}
\Needspace{5\baselineskip}
\pdfbookmark[2]{Werthe von Leibrenten und Lebensversicherungen für Männer.}{witwen-maenner}
\begin{center}\textit{Werthe von Leibrenten und Lebensversicherungen für Männer.}\end{center}
\noindent\hspace*{1.2em} Die Tafel für die Leibrenten der Männer hat \textsc{Gauss} zu einer genauen Berechnung des Einflusses benutzt, den diejenige Bestimmung der Statuten, dass ein Wiederaustritt des lebenden Mitgliedes nicht gestattet sein solle, haben würde. Er findet, dass für die 42 verheiratheten Mitglieder der Bilanzrechnung vom 1. Oct. 1845 der Zeitwerth der Beiträge sich dadurch um 898 Thl. bei \(3\tfrac{1}{2}\) Proc. und um 665 Thl. bei 4 Proc. vermehren würde. Ausserdem hat er mit Hülfe dieser Tafel einige Rechnungen über die Werthe von Lebensversicherungen ausgeführt und dabei für ein jetziges Mannesalter von \(M\) Jahren \(\rho-(1-\rho)\Phi M\) als Zeitwerth der am Ende des Todesjahres auszuzahlenden Münzeinheit genommen.]

\Needspace{6\baselineskip}
\pdfbookmark[2]{Werth der bestehenden Witwenpension Pm.}{witwen-pension-pm}
\begin{center}\textit{Werth der bestehenden Witwenpension \(Pm\).}\end{center}
\noindent \(m\)\quad jetziges Alter der Witwe.\par
\noindent \(\varphi m\)\quad Werth der Pension wenn jährlich und nur an noch Lebende gezahlt wird.\par
\noindent \(\rho\)\quad Discontofactor (\(=\tfrac{25}{26}\) für 4 Proc., \(=\tfrac{200}{207}\) für \(3\tfrac{1}{2}\) Proc.) oder der Zinsfuss so verstanden, dass \(\rho\) nach einem Jahre auf 1 anwächst.\par
\noindent \(fx\)\quad Lebende des Alters \(x\) nach Angabe der Mortalitätstafel.

[Die Bilanz wird für den 1. October eines bestimmten Jahres berechnet und dieser Zeitpunkt hier überall nur kurz der jetzige genannt. Nach dem Regulative vom 11. October 1833 und den später ergangenen Verfügungen die Professoren Witwenkasse betreffend wird die Pension in halbjährigen am 1. April und am 1. October jeden Jahres fälligen Raten ausbezahlt und erlischt bei Witwen mit dem Sterbemonate, welcher zu voll bezahlt wird. Mit Rücksicht hierauf ist für den wahrscheinlichen Jetztwerth \(Pm\) einer mit der Münzeinheit jährlich auszuzahlenden Witwenpension:]
\[
\begin{aligned}
fm\mathord{.}Pm&=\tfrac{1}{12}\rho^{\frac{1}{2}}\{fm+f(m+\tfrac{1}{12})+f(m+\tfrac{2}{12})+.\ .+f(m+\tfrac{5}{12})\}\\
&\quad+\tfrac{1}{12}\rho\{f(m+\tfrac{6}{12})+f(m+\tfrac{7}{12})+.\ .\ .\ .+f(m+\tfrac{11}{12})\}\\
&\quad+\text{u.\,s.\,w.}
\end{aligned}
\]
\noindent oder
\[
fm\mathord{.}Pm=\tfrac{1}{2}\rho^{\frac{1}{2}}f(m+\tfrac{5}{24})+\tfrac{1}{2}\rho f(m+\tfrac{17}{24})+\tfrac{1}{2}\rho^{\frac{3}{2}}f(m+\tfrac{29}{24})+\text{u.\,s.\,w.}
\]
\noindent wofür genommen werden kann
\[
=\rho^{\frac{3}{4}}f(m+\tfrac{11}{24})+\rho^{\frac{7}{4}}f(m+\tfrac{35}{24})+\text{u.\,s.\,w.}
\]
\noindent Die Tafel gibt
\[
fm\mathord{.}\varphi m=\rho f(m+1)+\rho\rho f(m+2)+.\ .\ .
\]
\noindent also
\[
\begin{aligned}
Pm&=\rho^{-\frac{1}{4}}\mathord{.}\frac{f(m-\tfrac{13}{24})}{fm}\mathord{.}\varphi(m-\tfrac{13}{24})\\
&=\rho^{-\frac{1}{4}}\mathord{.}g(m+\tfrac{11}{48})^{\frac{13}{24}}\mathord{.}\varphi(m-\tfrac{13}{24})
\end{aligned}
\]
\relax
[wenn man die in obiger Tafel unter decr. in Einheiten der siebenten Decimale enthaltenen Werthe von \(\log f(m-1)-\log fm\) mit \(\log\mathord{.}gm\) bezeichnet. Eine zur Berechnung von Hülfstafeln etwas bequemere Formel entsteht, wenn man bei jeder ganzen Zahl \(n\) und jedem echten Bruche \(k\) für \(f(n+k)\) die Grösse \((k-1)fn+kf(n+1)\) setzt, nemlich:]
